% ==============================================================
%  Chapter 1 -- Introduction and State of the Art
% ==============================================================

\chapter{Introduction}
\label{ch:introduction}

% --------------------------------------------------------------
\section{Problem Definition}
% --------------------------------------------------------------

A \emph{Decision Support System} (DSS) is an interactive
computational tool designed to assist a decision-maker in the analysis
of unstructured problems by combining data, formal models and
analytical techniques~\cite{insua2005,howard2008foundations}. Within
this broad family, \emph{Influence
Diagrams}~(IDs)\cite{howard2005influence,koller2009pgm,shenoy2009eolss}
have become a particularly compact graphical formalism:
a single directed acyclic graph encodes simultaneously the
probabilistic structure of the problem (chance nodes), the
controllable choices of the agent (decision nodes) and her
preferences over outcomes (utility nodes). Unlike a fully expanded
decision tree, an ID allows the analyst to reason about the problem
at the level of variables and to compute the optimal policy by direct
evaluation of the graph using exact algorithms such as the one of
Shachter~\cite{shachter1986evaluating}.

The expressive power of an ID rests on two distinct ingredients.
The first is its \emph{qualitative structure} — the set of nodes and
arcs — which captures the conditional dependencies of the domain.
The second is its \emph{quantitative parameterisation}, consisting of
the conditional probability tables (CPTs) of all chance nodes and the
numerical entries of the utility function. While the structure can
often be elicited from a domain expert in a relatively small number of
modelling sessions, the parameterisation grows exponentially with the
connectivity of the graph and becomes the real bottleneck of any
realistic ID-based DSS~\cite{bielza2010modeling}. A well-known
illustration is the neonatal-jaundice ID of G{\'o}mez et
al.~\cite{gomez2007jaundice}: more than 8\,000 probability entries
and over a hundred utility values had to be elicited from clinicians
before the model could be deployed, and the very same effort is
documented as the dominant cost of the project in the companion
analysis of Bielza et al.~\cite{bielza2000issues}.

This thesis addresses precisely this parameter-elicitation problem.
We assume that the structure of the ID is given, that all variables
are discrete and that the model contains a single value node. Under
these assumptions, the question we tackle can be stated as follows:

\begin{quote}\itshape
Given the qualitative structure of an Influence Diagram and a small
set of decision rules provided by a domain expert, can we
automatically recover a parameter vector — CPTs and utilities — that
reproduces those rules, thereby relieving the expert from the burden
of specifying every numerical entry of the model?
\end{quote}

% --------------------------------------------------------------
\section{Motivation and Importance}
% --------------------------------------------------------------

Several factors make this problem both timely and important.

\paragraph{The elicitation bottleneck.} Direct elicitation techniques
— reference lotteries, indifference judgements, pairwise comparisons —
remain the de-facto standard for the construction of probabilistic
decision models~\cite{bielza2010modeling}. They demand long
interviews with scarce experts, are cognitively demanding for the
interviewee and are well known to introduce inconsistencies that
require subsequent revision. In domains such as clinical medicine,
public-health planning or safety-critical engineering, the cost of
expert time is the dominant factor preventing the wider adoption of
ID-based decision tools.

\paragraph{Scarcity of historical data.} A natural alternative to
direct elicitation is to learn the parameters of the model from
historical data via maximum-likelihood or expectation–maximisation
estimators, as pioneered by Ezawa and Gupta in the context of
IDs~\cite{ezawa1997learning}. However, the regime in which most
clinical and engineering problems live is precisely the opposite of
what such statistical methods require: there are very few annotated
historical cases, but a domain expert is usually able to verbalise a
small number of high-level decision rules of the form
``\textit{if condition $\mathbf{x}$ holds, then take action $a$}''.

\paragraph{Decisions instead of probabilities.} Eliciting rules is
cognitively much closer to clinical or managerial practice than
eliciting numerical probabilities. Practitioners reason naturally in
terms of recommended actions under prototypical scenarios, and reach
agreement on rules far more easily than on probability tables.
Exploiting rules as a supervisory signal therefore aligns the
elicitation process with the way experts actually think.

\paragraph{Use cases.} Concrete domains where the proposed approach
is of immediate practical interest include:
\begin{itemize}
    \item \textbf{Clinical decision support.}
          Risk stratification, triage, treatment selection or the
          management of chronic diseases — the neonatal-jaundice ID
          of~\cite{gomez2007jaundice} is a canonical example, and
          the demand for evidence-based decision aids in clinical
          practice is by now well documented~\cite{dowding2004evidence}.
    \item \textbf{Industrial maintenance and safety.}
          Inspection scheduling, replacement-vs-repair choices, and
          safety-critical operations in which expert rules are
          available but historical incident data are sparse.
    \item \textbf{Public policy and emergency response.}
          Vaccination campaigns, evacuation plans and resource
          allocation under uncertainty, where domain rules exist but
          full probabilistic models do not.
    \item \textbf{Cyber-security and fraud detection.}
          Operational analysts maintain rule-based playbooks that can
          be turned into the supervisory signal of an ID-based
          decision module.
\end{itemize}

% --------------------------------------------------------------
\section{Scope}
% --------------------------------------------------------------

The field of decision support systems is vast. To keep the
contribution focused, the work in this thesis is restricted to:
\begin{enumerate}
    \item Influence Diagrams with \emph{discrete} chance and decision
          variables;
    \item a \emph{single} value (utility) node;
    \item a \emph{pre-specified} dependency graph (the structure is
          provided as an input and is not subject to learning);
    \item supervision signals consisting of a finite set of
          \emph{partial decision rules} produced by a domain expert.
\end{enumerate}
Extensions to dynamic IDs, multi-agent settings (MAIDs) and
continuous variables are discussed only in the perspectives chapter.

% --------------------------------------------------------------
\section{State of the Art}
\label{sec:state-of-the-art}
% --------------------------------------------------------------

The problem of obtaining the numerical parameters of a graphical
decision model has been approached from several angles. We organise
the literature along four complementary lines and then position the
present work within them.

\subsection{Statistical learning from data (MLE / EM)}

The classical route to parameterising any graphical model is
maximum-likelihood estimation. When the data are fully observed the
CPTs reduce to relative frequencies; with hidden variables or missing
entries the Expectation–Maximisation algorithm is the standard
choice. In the specific context of IDs, the seminal work of Ezawa
and Gupta~\cite{ezawa1997learning} combined EM for the CPTs with a
greedy search for the causal structure, and extended the procedure to
the utility node via regression on observed outcomes. The approach
is implemented in essentially every graphical-models toolkit (GeNIe,
pgmpy, bnlearn) but requires a sizeable corpus of labelled historical
cases, which in many real applications is simply not available.

\subsection{Direct expert elicitation}

When historical data are scarce, the established alternative is the
structured elicitation of probabilities and utilities through
techniques such as reference lotteries, indifference judgements and
canonical models (Noisy-OR, Noisy-AND, additive utilities)
\cite{bielza2010modeling}. The literature on multi-criteria decision
making, including the doctoral work of Shahzad on the integration of
Bayesian networks and MCDM~\cite{shahzad2022mcdm}, has produced a
mature catalogue of methods to reduce the cognitive load on the
expert. The bottleneck nevertheless persists: every CPT row and
every utility entry has to be confirmed by a human, and the size of
the table grows exponentially with the number of parents.

\subsection{Evolutionary algorithms and metaheuristics}

A third line of work treats parameter elicitation as an
\emph{inverse optimisation} problem: rather than asking the expert
for every parameter, one searches the parameter space for a
configuration that reproduces a weak supervisory signal. Genetic
Algorithms have been used inside the UPM CIG group for this very
purpose~\cite{henriquez2024tfm}, and swarm-intelligence techniques
(PSO, ACO) have been applied to the related problem of learning
unknown behaviours in Interactive Dynamic IDs by Pan et
al.~\cite{pan2025iddids}. Within evolutionary computation, the
family of \emph{Estimation of Distribution Algorithms}
(EDAs)~\cite{larranaga2002eda} has matured into a powerful framework
that, instead of recombining individuals via crossover and mutation,
explicitly learns a probabilistic model of the promising regions of
the search space at every generation. Recent surveys
\cite{larranaga2024edas} report consistent improvements over
classical GAs in continuous, high-dimensional and constrained
problems, and the open-source library
\textsc{EDAspy}~\cite{soloviev2024edaspy} now provides
production-quality implementations of UMDA, EGNA, KEDA and SPEDA, as
well as the semiparametric variant SPEDA~\cite{soloviev2023speda}.

\subsection{Generative and LLM-based elicitation}

The most recent trend exploits large language models as
``in-silico'' experts. Prompted with the structure of the network
and natural-language descriptions of the variables, models such as
GPT-4o, Gemini Pro and Claude 3.5 have been shown to produce CPTs
whose Kullback–Leibler divergence with respect to expert-elicited
tables is comparable to the noise observed across different human
experts~\cite{llm4bn2025}. Closely related, Pan et al.~use
Variational Auto-encoders to generate plausible models of the
unobserved agents in Interactive Dynamic IDs~\cite{pan2024vae}.
These methods are still in their infancy but already suggest powerful
hybrid pipelines in which a generative model produces a warm-start
configuration that is subsequently refined by a metaheuristic.

\subsection{Exact optimisation and reinforcement learning}

For completeness, two further lines must be mentioned. The
\emph{Decision Programming} framework of Salo et
al.~\cite{salo2022decisionprog}, refined by Hankimaa et
al.~\cite{hankimaa2024mip}, recasts the ID as a mixed-integer linear
program that can be solved by off-the-shelf optimisation engines.
G{\'o}mez Castellanos and Terashima-Mar{\'\i}n~\cite{gomez2024rl} use
deep reinforcement learning (PPO) to solve Mixed IDs by treating the
optimal policy as the output of an actor-critic. Both lines tackle
the dual problem to the one studied here: they fix the parameters and
optimise the policy, whereas in this thesis the policy
(the expert rules) is fixed and the parameters are searched.
The recent survey of Vonk and
Gonzalez-Soto~\cite{vonk2026graphical} gives the most up-to-date
panorama of these complementary approaches.

\subsection{The inverse problem: rules from CPTs}

A complementary line of work, also rooted at UPM, attacks the
\emph{inverse} of the problem studied in this thesis. Fern{\'a}ndez
del Pozo, Bielza and Lucas~\cite{fdezdelpozo2010cluster} propose the
\textsc{Cluster-KBM2L} methodology, which starts from a fully
specified CPT (more than 8\,000 entries in the case of their
non-Hodgkin lymphoma ID) and extracts \emph{symbolic prognostic
profiles} — qualitative rules of the form ``high-risk patient with
favourable chemotherapy response'' — by clustering the conditional
distributions and synthesising interpretable descriptions for each
cluster. Both works thus address the same interface between the
numerical and the symbolic representation of an ID, but from
opposite directions: \cite{fdezdelpozo2010cluster} goes from CPTs
to rules, whereas this thesis goes from rules to CPTs.

\subsection{Research gap and positioning of this work}

The systematic review summarised above leads to a clear conclusion:
no published work combines \emph{modern EDAs} with the
\emph{reconstruction of ID parameters from a small set of expert
decision rules}. The closest precedent is the genetic-algorithm
approach of Henr{\'\i}quez~\cite{henriquez2024tfm}. The contribution
of the present thesis is to replace the GA with an EDA that
\emph{explicitly} models the dependency structure of the parameter
space, and to study empirically the impact of this richer
probabilistic model on convergence and final accuracy.

% --------------------------------------------------------------
\section{Objectives and Contributions}
% --------------------------------------------------------------

The overall objective of this thesis is to develop and validate a
proof-of-concept system that estimates the quantitative parameters of
a regular Influence Diagram from a partial expert rule base, using
Estimation of Distribution Algorithms as the search engine. To
achieve this overarching goal we set the following specific
objectives:
\begin{enumerate}
    \item Formalise the parameter-recovery problem as a constrained
          black-box optimisation over the joint space of CPTs and
          utilities.
    \item Design an encoding of the parameter vector compatible with
          continuous EDAs and capable of guaranteeing the validity of
          every decoded CPT.
    \item Implement the resulting system on top of the
          \textsc{EDAspy} library~\cite{soloviev2024edaspy} and the
          \textsc{SMILE} inference engine~\cite{genie}, using the
          neonatal-jaundice ID of~\cite{gomez2007jaundice} as the
          benchmark.
    \item Empirically evaluate the impact of the main
          hyperparameters (population size, selection ratio, dead
          iterations) on the quality of the recovered parameters and
          on the run-time of the algorithm.
    \item Discuss the limits of the proposed approach and outline a
          set of natural extensions (LLM-informed initialisation,
          canonical-model parameterisation, SPEDA, parallel EDAs).
\end{enumerate}

% --------------------------------------------------------------
\section{Thesis Outline}
% --------------------------------------------------------------

The remainder of the document is organised in four further chapters.

\begin{itemize}
    \item \textbf{Chapter~\ref{ch:framework}} develops the
          theoretical framework: Decision Support Systems, the
          formal definition of an Influence Diagram and its
          evaluation, the Knowledge-Acquisition problem that
          motivates this thesis, and a self-contained presentation
          of Estimation of Distribution Algorithms with the variants
          relevant for our setting.
    \item \textbf{Chapter~\ref{ch:implementation}} describes the
          implementation of the proof of concept: the architecture
          of the system, the encoding and decoding of the parameter
          vector, the design of the fitness function and the
          configuration of the EDA loop with \textsc{EDAspy}.
    \item \textbf{Chapter~4} reports the experimental study: the
          benchmark ID, the protocol of the hyperparameter sweep,
          the comparison between the rules generated by the
          optimised model and those provided by the expert, and a
          discussion of the empirical results.
    \item \textbf{Chapter~5} concludes the work, summarising the
          contributions and outlining the research directions
          identified in the state of the art that constitute
          natural lines of future work.
\end{itemize}
